\documentclass[a4paper, 12pt]{article}
\usepackage[utf8]{inputenc}
\usepackage[T1]{fontenc}
\usepackage[french]{babel}
\usepackage{graphicx}
\usepackage{amsmath}
\usepackage{amssymb}
\usepackage{hyperref}
\usepackage{vmargin}
\setmargnohfrb{30mm}{30mm}{30mm}{30mm} %ltrb

\pagestyle{plain}

\title{Projet Dev Web II}
\author{Par : \textsc{Defay} Adrien, \textsc{Breuil} Dorian et \textsc{Peyrard} Gaultier}
\date{\today}

\begin{document}

\maketitle

\newpage

\tableofcontents
\setcounter{tocdepth}{2}  

\newpage 

\section{Presentation du Projet}

Le projet présenté est un "Google drive", c'est à dire un document texte que plusieurs utilisateurs peuvent modifier en même temps.

Le projet suit une structure web java et possède 4 dossiers principaux : 
\begin{itemize}
	\item Controller ;
	\item Dao ;
	\item Model ;
	\item et View.
\end{itemize}

Controller permet de gérer toutes les servlets, Dao contient tout ce qui est en rapport avec la base de données, Model contient toutes les objets java et View contient tous les fichier jsp et html.
Quand au css, on utilise Bulma, un dossier java script dans Webapp et la base de donnée est à la racine du projet


\end{document}